\documentclass{article}

\usepackage{amsmath,amssymb}
\usepackage{xurl}
\usepackage[colorlinks=true,linkcolor=blue,citecolor=blue,urlcolor=blue]{hyperref}
\setlength{\emergencystretch}{2em}

% Shared commands for notes in docs/paper-gaps/.
% Notation follows the LDT paper (references/ldt-paper/).

\newcommand{\Lean}{\textsc{Lean}}
\newcommand{\leanid}[1]{\textsf{#1}}
\newcommand{\ghissue}[1]{\href{https://github.com/LionSR/MIPStarRE/issues/#1}{GitHub issue \##1}}
\newcommand{\ghpr}[1]{\href{https://github.com/LionSR/MIPStarRE/pull/#1}{GitHub PR \##1}}

% --- Paper-compatible notation ---
\newcommand{\F}{\mathbb{F}}
\newcommand{\N}{\mathbb{N}}
\newcommand{\C}{\mathbb{C}}
\newcommand{\tr}{\operatorname{tr}}
\newcommand{\ot}{\otimes}
\newcommand{\E}{\mathop{\bf E\/}}
\newcommand{\bx}{\mathbf{x}}
\newcommand{\by}{\mathbf{y}}
\newcommand{\bu}{\mathbf{u}}
\newcommand{\bv}{\mathbf{v}}

% Approximation relation (paper uses \approx_\eps and \simeq_\zeta)
\newcommand{\approxeps}{\approx_{\varepsilon}}
\newcommand{\simgeq}{\simeq_{\zeta}}

% Polynomial spaces (paper uses \polyfunc, \polysub, \polymeas)
\newcommand{\polyfunc}[3]{\operatorname{Poly}(#1,#2,#3)}
\newcommand{\polysub}[3]{\operatorname{PolySub}(#1,#2,#3)}
\newcommand{\polymeas}[3]{\operatorname{PolyMeas}(#1,#2,#3)}


\title{Model Note: Comparing a Cited Argument with a Formal Statement}
\author{}
\date{\today}

\begin{document}
\maketitle

% This file is a model for notes in docs/paper-gaps/. See policy.tex for the
% full style guide. Key rules:
%   - Notation must match the cited source (use command.tex macros).
%   - Lean identifiers, file paths, and issue numbers go in footnotes only.
%     Use \leanid{Name} for Lean identifiers, \ghissue{N} for issues,
%     \ghpr{N} for pull requests, \doclink{path/to/File.lean} for doc-gen links.
%   - Write in the first-person plural, active voice, direct sentences.
%   - Begin with a concise estimate when the note concerns formalization work.
%   - Put Key theorem forms as Section 2, immediately after At a glance.
%   - No \noindent\textbf{Verdict:} labels: the Conclusion section is the verdict.
%   - Keep blueprint/formal comparison and conclusion compact.
%   - The body must read as mathematical prose, not as an audit log.
%   - Compile from the repository root with pdflatex.

\section{At a glance}

\begin{itemize}
  \item \textbf{Difficulty.} Easy, medium, or hard, with one sentence explaining why.
  \item \textbf{Estimated weight.} Roughly how many new definitions or theorem
  interfaces and how many lines of formal proof are expected.  Give a range.
  \item \textbf{Mathlib/project split.} Estimate the percentage that should be
  direct Mathlib reuse versus project-local development.
  \item \textbf{Key Mathlib inputs.} List the main mathematical APIs or modules,
  not every imported file.
\end{itemize}

\section{The source assertion}

We examine a step in \cite{Ji2020LowIndividualDegree}.%
\footnote{Tracked in \ghissue{NNN}. The paper passage is at
\path{references/ldt-paper/FILE.tex}, lines~A--B.}
In the notation of the paper, the cited source needs
\[
  \mathcal H \Longrightarrow C.
  \tag{\texttt{source-label}}
\]
Here $\mathcal H$ is the available hypothesis set, and $C$ is the conclusion
used later.

\section{Key theorem forms}

The formalization should isolate the following mathematical statements.
\begin{enumerate}
  \item \textbf{Source theorem form.}  State the theorem exactly as the paper uses it.
  \item \textbf{Adapter theorem form.}  State any translation between the paper
  objects and the formal objects.
  \item \textbf{Downstream theorem form.}  State the exact conclusion consumed
  later in the proof.
\end{enumerate}
Keep declaration names in a footnote if they are needed for traceability.%
\footnote{For example, the formal declarations might be \leanid{modelSource},
\leanid{modelAdapter}, and \leanid{modelDownstream} in
\doclink{MIPStarRE/Path/To/File.lean}.}

\section{Comparison with the blueprint and formalization}

The blueprint should either match the cited source or explicitly record the
changed theorem form.  The formal development proves the corrected or adapted
statement
\[
  \mathcal H \land \mathcal K \Longrightarrow C,
\]
where $\mathcal K$ is the additional input isolated above.  Explain in prose
whether $\mathcal K$ is already available, still open, or unnecessary after a
different argument.

\section{Conclusion}

State the verdict plainly.  For example: the source statement is correct and the
remaining work is formalization; the source proof needs an additional
hypothesis; or there is no discrepancy and no further action is needed.

\bibliographystyle{alpha}
\bibliography{references}

\end{document}
